\section{Methodology and Pipeline}

Methodologically, the report keeps a functional connectivity perspective (not effective connectivity) and emphasizes robustness against volume conduction and false positives \citep{Cohen2014,Vinck2011WPLI,Cao2022EEGConnectivity}. The working order of the pipeline is intentional: data ingestion and quality checks, deterministic filtering, ICA decomposition and cleaning, segment-level normalization and artifact rejection, optional CSD as a planned spatial improvement, band-limited wPLI estimation, and finally surrogate-based statistical inference. This order is preserved because each stage prepares a specific assumption required by the next one, so if a QC check fails the rerun should start from the corresponding local phase.

From an internal development standpoint, the workflow is implemented as modular phases with explicit persistence between steps. Intermediate outputs are serialized so that any phase can be rerun without recomputing the whole chain, which is essential for debugging and sensitivity analysis. In practice, this also allows quick comparisons between parameter sets (for example, artifact thresholds, segment length, or surrogate count) while keeping full traceability of generated artifacts \citep{Blinowska-Cieslak2022,Sanei2022}.
The operational route can be read as three contiguous blocks: acquisition harmonization and cleaning (Phases 0--2), segment-level preparation and signal conditioning (Phases 3--6), and connectivity plus inferential validation (Phases 7--8). This grouping is useful when debugging because it reduces the search space: failures before ICA are usually data/metadata issues, failures in the middle block are usually segmentation or thresholding decisions, and failures at the end typically involve band definitions, connectivity assumptions, or surrogate settings.
\sidenote{%
\footnotesize
\textbf{Pipeline map}\\
P0: \texttt{BIDS.jl} $\rightarrow$ \texttt{src/Setup/IO.jl}\\
P1: \texttt{Preprocessing.jl}\\
P2a--2b: \texttt{ICA.jl}, \texttt{ICA\_cleaning.jl}\\
P3--6: \texttt{Processing.jl}\\
P7: \texttt{Connectivity.jl}\\
P8: \texttt{Surrogate.jl} (documented/planned)\\
\textit{If one phase fails, rerun locally.}
}

\subsection{Pipeline summary and execution order}

The implementation is organized as a sequential chain with local persistence between phases. This allows deterministic reruns from the failure point without recomputing all previous stages.

\begin{tcolorbox}[
  enhanced,
  breakable,
  colback=blue!2,
  colframe=blue!55!black,
  coltitle=black,
  boxrule=0.6pt,
  arc=2mm,
  left=1.2mm,
  right=1.2mm,
  top=1mm,
  bottom=1mm,
  title=\textbf{Pipeline phases and execution order}
]
\begin{enumerate}[leftmargin=*,itemsep=2pt,topsep=2pt]
\item \textbf{Phase 0 (IO/BIDS):} load raw EEG, validate metadata consistency, and compute initial QC indicators.
\item \textbf{Phase 1 (Filtering):} apply deterministic filtering in a fixed order (notch 50\,Hz, bandreject 100\,Hz, high-pass 0.5\,Hz, low-pass 150\,Hz).
\item \textbf{Phases 2a--2b (ICA):} decompose into ICs and reconstruct clean EEG after artifact-component removal.
\item \textbf{Phases 3--6 (Segment processing):} segment, baseline-correct, reject artifacts, and produce quality-controlled tensors.
\item \textbf{Phase 7 (Connectivity):} apply CSD and estimate band-limited wPLI matrices.
\item \textbf{Phase 8 (Inference):} build surrogate null distributions, compute p-values, and apply multiple-comparison control.
\end{enumerate}
\end{tcolorbox}

The current repository separates executable source modules from narrative Pluto notebooks. Table~\ref{tab:pipeline-traceability} provides a compact traceability map between the published documentation layer and the code artifacts that actually run each stage.

\begin{fullwidth}
\begin{table}[H]
\centering
\scriptsize
\caption{Traceability map between pipeline phase, documentation notebook, executable source, and main persisted outputs.}
\label{tab:pipeline-traceability}
\begin{adjustbox}{max width=0.98\fulltextwidth}
\begin{tabular}{@{}>{\raggedright\arraybackslash}p{0.11\fulltextwidth}>{\raggedright\arraybackslash}p{0.15\fulltextwidth}>{\raggedright\arraybackslash}p{0.29\fulltextwidth}>{\raggedright\arraybackslash}p{0.32\fulltextwidth}>{\raggedright\arraybackslash}p{0.08\fulltextwidth}@{}}
\toprule
\textbf{Phase} & \textbf{Pluto} & \textbf{Executable source} & \textbf{Persisted outputs} & \textbf{Status} \\
\midrule
BIDS/IO & \texttt{BIDS.jl}; \texttt{Project\_Setup.jl} & \texttt{src/BIDS/*.jl}; \texttt{src/Setup/IO.jl} & \bids{data/BIDS/}; \bids{data/Preprocessing/IO/dict\_EEG.bin} & Active \\
\addlinespace
Filtering & \texttt{Preprocessing.jl} & \texttt{src/Preprocessing/filtering.jl} & \bids{data/Preprocessing/filtering/}; filter-response and PSD figures & Active \\
\addlinespace
ICA & \texttt{ICA.jl} & \texttt{src/ICA/ICA.jl}; \texttt{src/ICA/ICA\_cleaning.jl} & \bids{data/Processing/ICA/}; \bids{results/Processing/figures/ICA\_cleaning/}; IC audit tables & Active \\
\addlinespace
Processing & \texttt{Processing.jl} & \texttt{src/Processing/segmentation.jl}; \texttt{baseline*.jl}; \texttt{artifact\_rejection.jl} & Segmented tensors, baseline-corrected tensors, rejected-epoch summaries & Active \\
\addlinespace
Spectral & \texttt{Spectral.jl} & \texttt{src/Spectral/FFT.jl} & \bids{data/Processing/FFT/}; band-power tables and spectral figures & Active \\
\addlinespace
Connectivity & \texttt{Connectivity.jl} & \texttt{src/Connectivity/CSD.jl}; \texttt{src/Connectivity/wPLI.jl} & \bids{data/Connectivity/CSD/}; \bids{data/Connectivity/wPLI/}; heatmaps and edge tables & Active \\
\addlinespace
Inference & \texttt{Surrogate.jl} & \texttt{src/Surrogate/Surrogate.jl} & Surrogate null matrices, empirical p-values, and FDR tables & Pending \\
\bottomrule
\end{tabular}
\end{adjustbox}
\end{table}
\end{fullwidth}

A compact step-to-script mapping is provided in the sidenote.
\sidenote{%
\scriptsize
\textbf{Pipeline mapping (Pluto $\rightarrow$ src)}\\
\textbf{Step 1} \texttt{Project\_Setup.jl}\\
\texttt{/BIDS/build\_eeg\_bids.jl}\\
\texttt{/BIDS/build\_participants.jl}\\
\texttt{/BIDS/validate\_bids.jl}\\[2pt]
\textbf{Step 2} \texttt{BIDS.jl}\\
\texttt{/Setup}; \texttt{/IO.jl}\\
\textbf{Step 3} \texttt{Preprocessing.jl}\\
\texttt{/Preprocessing/filtering.jl}\\
\textbf{Step 4} \texttt{ICA.jl}\\
\texttt{/ICA/ICA.jl}; \texttt{/ICA/ICA\_cleaning.jl}\\[2pt]
\textbf{Step 5} \texttt{Processing.jl}\\
\texttt{/Processing/segmentation.jl}\\
\texttt{/Processing/baseline.jl}\\
\texttt{/Processing/artifact\_rejection.jl}\\
\texttt{/Processing/baseline\_2st.jl}\\[2pt]
\textbf{Step 6} \texttt{Spectral.jl}\\
\texttt{/Spectral/FFT.jl}\\
\textbf{Step 7} \texttt{Connectivity.jl}\\
\texttt{/Connectivity/CSD.jl}; \texttt{/Connectivity/wPLI.jl}\\
\textbf{Step 8} \texttt{Surrogate.jl}\\
\texttt{/Surrogate/Surrogate.jl} (pending executable module)
}

\subsection{Phase 0 (IO/BIDS)}

Raw data are loaded and standardized into per-subject, per-condition structures with validated metadata. This stage is where channel labels, dimensions, and subject/condition consistency are enforced. It also generates initial quality statistics and exploratory plots used as a baseline before filtering.
\sidenote{%
\scriptsize
\textbf{Phase 0: graphical outputs (QC)}\\
Raw traces (single + multichannel)\\
Average PSD per channel\\
Band-power summaries\\
Channel statistics plots\\
Flags for suspicious channels\\[2pt]
\textbf{Stats (quick defs)}\\
$\mu$: mean,\; $\sigma$: std\\
RMS $= \sqrt{\frac{1}{N}\sum x_n^2}$\\
$\gamma_1$: skewness,\; $\gamma_2$: kurtosis
}
For each channel signal $x=\{x_n\}_{n=1}^{N}$ in this first phase, descriptive statistics are computed to characterize baseline behavior and detect suspicious channels before filtering. The mean is used to track DC offset and the standard deviation to summarize dispersion.
\sidenote{%
\scriptsize
\textbf{Mean}\\
\textit{Average signal level (DC offset estimate).}\\
$\mu=\frac{1}{N}\sum_{n=1}^{N}x_n$
}
\sidenote{%
\scriptsize
\textbf{Standard deviation}\\
\textit{Dispersion of samples around the mean.}\\
$\sigma=\sqrt{\frac{1}{N-1}\sum_{n=1}^{N}(x_n-\mu)^2}$
}

Extremes are monitored with minimum and maximum values to detect clipping or saturation, and RMS is used as a compact measure of signal energy amplitude.
\sidenote{%
\scriptsize
\textbf{Minimum / maximum}\\
\textit{Observed amplitude bounds in the channel.}\\
$x_{\min}=\min(x)$\\
$x_{\max}=\max(x)$
}
\sidenote{%
\scriptsize
\textbf{RMS}\\
\textit{Root-mean-square amplitude (signal energy proxy).}\\
$\mathrm{RMS}=\sqrt{\frac{1}{N}\sum_{n=1}^{N}x_n^2}$
}

Shape descriptors are also tracked to characterize asymmetry and impulsive behavior in channel distributions.
\sidenote{%
\scriptsize
\textbf{Skewness}\\
\textit{Asymmetry of the amplitude distribution.}\\
$\gamma_1=\frac{1}{N}\sum\left(\frac{x_n-\mu}{\sigma}\right)^3$
}
\sidenote{%
\scriptsize
\textbf{Kurtosis}\\
\textit{Tail heaviness / impulsiveness of the distribution.}\\
$\gamma_2=\frac{1}{N}\sum\left(\frac{x_n-\mu}{\sigma}\right)^4$
}
\begin{fullwidth}
\begin{figure}[H]
\centering
\resizebox{0.92\textwidth}{!}{%
\begin{tikzpicture}[x=1cm,y=1cm,>=Latex,font=\scriptsize]
  % ===== Panel 1: Mean =====
  \begin{scope}[shift={(0,2.9)}]
    \draw[black!65] (0,0) rectangle (3.9,2.1);
    \node[anchor=west,font=\scriptsize\bfseries] at (0.08,1.92) {Mean};
    \draw[-Latex,black!55] (0.45,0.35)--(3.45,0.35);
    \draw[-Latex,black!55] (0.45,0.35)--(0.45,1.70);
    \node[black!65,anchor=west] at (3.00,0.42) {$t$};
    \node[black!65,rotate=90,anchor=south] at (0.18,0.52) {$\mu$V};
    \draw[black!45] (0.35,0.95)--(3.55,0.95);
    \draw[blue!70!black,thick]
      (0.35,0.75) .. controls (0.9,1.45) and (1.5,0.45) .. (2.0,1.05)
      .. controls (2.45,1.35) and (3.0,0.7) .. (3.55,1.15);
    \draw[red!75!black,dashed,thick] (0.45,1.05)--(3.55,1.05);
    \node[red!75!black,anchor=west] at (3.00,1.13) {$\mu$};
  \end{scope}

  % ===== Panel 2: Standard deviation =====
  \begin{scope}[shift={(4.2,2.9)}]
    \draw[black!65] (0,0) rectangle (3.9,2.1);
    \node[anchor=west,font=\scriptsize\bfseries] at (0.08,1.92) {Std.\ deviation};
    \draw[-Latex,black!55] (0.45,0.35)--(3.45,0.35);
    \draw[-Latex,black!55] (0.45,0.35)--(0.45,1.70);
    \node[black!65,anchor=west] at (3.00,0.42) {$t$};
    \node[black!65,rotate=90,anchor=south] at (0.18,0.52) {$\mu$V};
    \draw[black!45] (0.4,0.45)--(0.4,1.75);
    \draw[black!45] (3.5,0.45)--(3.5,1.75);
    \draw[blue!70!black,thick]
      (0.4,0.6) .. controls (1.35,1.65) and (2.55,1.65) .. (3.5,0.6);
    \draw[red!75!black,dashed,thick] (1.95,0.45)--(1.95,1.75);
    \draw[black!70,<->] (1.95,1.25)--(2.85,1.25);
    \node[anchor=west] at (2.9,1.25) {$\sigma$};
  \end{scope}

  % ===== Panel 3: Min / Max =====
  \begin{scope}[shift={(8.4,2.9)}]
    \draw[black!65] (0,0) rectangle (3.9,2.1);
    \node[anchor=west,font=\scriptsize\bfseries] at (0.08,1.92) {Min / Max};
    \draw[-Latex,black!55] (0.45,0.35)--(3.45,0.35);
    \draw[-Latex,black!55] (0.45,0.35)--(0.45,1.70);
    \node[black!65,anchor=west] at (3.00,0.42) {$t$};
    \node[black!65,rotate=90,anchor=south] at (0.18,0.52) {$\mu$V};
    \draw[black!45] (0.35,0.95)--(3.55,0.95);
    \draw[blue!70!black,thick]
      (0.35,0.95)--(0.9,1.55)--(1.35,0.55)--(1.9,1.1)--(2.45,0.35)--(3.0,1.35)--(3.55,0.8);
    \fill[red!80!black] (2.45,0.35) circle (1.4pt);
    \fill[green!60!black] (0.9,1.55) circle (1.4pt);
    \node[red!80!black,anchor=west] at (2.52,0.38) {$x_{\min}$};
    \node[green!50!black,anchor=west] at (0.97,1.58) {$x_{\max}$};
  \end{scope}

  % ===== Panel 4: RMS =====
  \begin{scope}[shift={(0,0)}]
    \draw[black!65] (0,0) rectangle (3.9,2.1);
    \node[anchor=west,font=\scriptsize\bfseries] at (0.08,1.92) {RMS};
    \draw[-Latex,black!55] (0.45,0.35)--(3.45,0.35);
    \draw[-Latex,black!55] (0.45,0.35)--(0.45,1.70);
    \node[black!65,anchor=west] at (3.00,0.42) {$t$};
    \node[black!65,rotate=90,anchor=south] at (0.18,0.52) {$\mu$V};
    \draw[black!45] (0.35,0.95)--(3.55,0.95);
    \draw[blue!70!black,thick]
      (0.35,1.35).. controls (0.9,0.5) and (1.35,1.55) .. (1.9,0.8)
      .. controls (2.35,0.45) and (2.95,1.45) .. (3.55,0.7);
    \draw[orange!85!black,dashed,thick] (0.35,0.35)--(3.55,0.35);
    \node[orange!85!black,anchor=west] at (2.45,0.45) {RMS};
  \end{scope}

  % ===== Panel 5: Skewness =====
  \begin{scope}[shift={(4.2,0)}]
    \draw[black!65] (0,0) rectangle (3.9,2.1);
    \node[anchor=west,font=\scriptsize\bfseries] at (0.08,1.92) {Skewness};
    \draw[-Latex,black!55] (0.45,0.35)--(3.45,0.35);
    \draw[-Latex,black!55] (0.45,0.35)--(0.45,1.70);
    \node[black!65,anchor=west] at (3.00,0.42) {$t$};
    \node[black!65,rotate=90,anchor=south] at (0.18,0.52) {$\mu$V};
    \draw[black!45] (0.4,0.45)--(0.4,1.75);
    \draw[black!45] (3.5,0.45)--(3.5,1.75);
    \draw[blue!70!black,thick]
      (0.4,0.6) .. controls (1.0,1.6) and (1.9,1.45) .. (2.5,0.95)
      .. controls (3.0,0.65) and (3.25,0.58) .. (3.5,0.55);
    \node[anchor=west] at (2.42,1.55) {right tail};
  \end{scope}

  % ===== Panel 6: Kurtosis =====
  \begin{scope}[shift={(8.4,0)}]
    \draw[black!65] (0,0) rectangle (3.9,2.1);
    \node[anchor=west,font=\scriptsize\bfseries] at (0.08,1.92) {Kurtosis};
    \draw[-Latex,black!55] (0.45,0.35)--(3.45,0.35);
    \draw[-Latex,black!55] (0.45,0.35)--(0.45,1.70);
    \node[black!65,anchor=west] at (3.00,0.42) {$t$};
    \node[black!65,rotate=90,anchor=south] at (0.18,0.52) {$\mu$V};
    \draw[black!45] (0.4,0.45)--(0.4,1.75);
    \draw[black!45] (3.5,0.45)--(3.5,1.75);
    \draw[blue!70!black,thick]
      (0.4,0.5) .. controls (1.35,1.9) and (2.55,1.9) .. (3.5,0.5);
    \draw[black!45,dashed] (0.4,0.75) .. controls (1.35,1.45) and (2.55,1.45) .. (3.5,0.75);
    \node[anchor=west] at (2.35,1.58) {high peak};
  \end{scope}
\end{tikzpicture}
}
\caption{Visual intuition for the basic statistics used in Phase 0 (mean, dispersion, extremes, RMS energy, asymmetry, and kurtosis).}
\end{figure}
\end{fullwidth}
In operational QC terms, channels with unusually low variance/std may indicate poor contact or flat recordings, while very high kurtosis or extreme min/max values can suggest transient artifacts.

\subsection{Phase 1 (Filtering)}

Preprocessing takes the output of Phase 0 (IO): the raw EEG signals organized as a channel-indexed dictionary or matrix, with temporal and channel metadata and optional quality statistics from the initial inspection. The goal is to remove line noise, hardware-related interference, and slow drift, and to limit bandwidth before band-specific connectivity analysis.

The implementation in \texttt{src/Preprocessing/filtering.jl} loads serialized data from \texttt{data/Preprocessing/IO/}, applies a cascade of digital filters, and at each step stores the result under \texttt{data/Preprocessing/filtering/}, plots the filter response (magnitude and phase), and compares average PSD before and after filtering.

All filters are Butterworth. High-pass and low-pass stages use \texttt{filtfilt} (zero-phase filtering), so effective order doubles and temporal alignment is preserved. The cascade is applied in this order:

\begin{enumerate}
\item \textbf{Notch (50\,Hz):} bandstop stage for power-line interference (Europe), design order 4 with 1\,Hz bandwidth around 50\,Hz.
\item \textbf{Bandreject (100\,Hz):} narrow bandstop stage (order 4, 1\,Hz bandwidth) to suppress possible hardware-related interference near 100\,Hz without broad attenuation of higher frequencies.
\item \textbf{High-pass (0.5\,Hz):} drift removal while preserving slow activity; design order 4 (effective 8 with \texttt{filtfilt}).
\item \textbf{Low-pass (150\,Hz):} upper-bandwidth limitation; design order 4 (effective 8 with \texttt{filtfilt}).
\end{enumerate}

Stage-by-stage outputs are stored to preserve traceability and to support isolated reruns of downstream ICA or spectral analyses.

\subsection{Phases 2a--2b (ICA)}

After preprocessing (Phase 1), ICA is applied to separate neural activity from artifactual sources (ocular, muscular, line-noise and abrupt transients). The decomposition follows the standard linear mixing model
\[
\mathbf{X}=\mathbf{A}\mathbf{S},\qquad \mathbf{S}=\mathbf{W}\mathbf{X},
\]
with $\mathbf{X}\in\mathbb{R}^{C\times N}$ (channels $\times$ samples), $\mathbf{A}$ the mixing matrix, and $\mathbf{W}$ the demixing matrix estimated by symmetric FastICA.
\sidenote{%
\footnotesize
\textbf{ICA implementation}\\
Phase 2a: \texttt{src/ICA/ICA.jl}\\
Phase 2b: \texttt{src/ICA/ICA\_cleaning.jl}\\
Inputs from \texttt{data/Preprocessing/filtering/}\\
Outputs to \texttt{data/Processing/ICA/}
}

In Phase 2a, the objective is to estimate independent components and their spatial maps from filtered EEG while preserving the full channel dimensionality ($k=C$). The input is the filtered matrix $\mathbf{X}\in\mathbb{R}^{C\times N}$ from Phase 1.

The first operation is centering, which removes per-channel mean offsets and yields $\mathbf{X}_c$. Then PCA whitening is applied so that transformed data have identity covariance and ICA can focus on higher-order structure:
\[
\mathbf{Z}=\mathbf{V}_{\text{whit}}\mathbf{X}_c,\qquad \mathbf{V}_{\text{whit}}=\mathbf{\Lambda}^{-1/2}\mathbf{E}^{\top}.
\]
Here, $\mathbf{E}$ contains eigenvectors of the covariance matrix and $\mathbf{\Lambda}$ the corresponding eigenvalues. After whitening, symmetric FastICA is iterated with $g(u)=\tanh(au)$ and symmetric decorrelation until convergence.

The output of this phase is the component matrix $\mathbf{S}$, the total demixing matrix $\mathbf{W}_{\text{total}}=\mathbf{W}\mathbf{V}_{\text{whit}}$, and the mixing matrix $\mathbf{A}=\mathbf{W}_{\text{total}}^{-1}$, which maps IC activity back to channel space.
\sidenote{%
\scriptsize
\textbf{FastICA outputs}\\
$\mathbf{S}$: independent components\\
$\mathbf{W}_{\text{total}}=\mathbf{W}\mathbf{V}_{\text{whit}}$\\
$\mathbf{A}=\mathbf{W}_{\text{total}}^{-1}$
}

\begin{algorithm}[H]
\caption{Operational pseudocode for symmetric FastICA (Phase 2a)}
\begin{algorithmic}
\Require $\mathbf{X}\in\mathbb{R}^{C\times N}$, $k=C$, $\texttt{max\_iter}$, $\tau$, $a$, seed
\State Center channels: $\mathbf{X}_c \gets \mathbf{X}-\mathrm{mean}(\mathbf{X},\text{dims}=2)$
\State PCA whitening: $\mathbf{Z}\gets \mathbf{\Lambda}^{-1/2}\mathbf{E}^{\top}\mathbf{X}_c$
\State Initialize $\mathbf{W}\in\mathbb{R}^{k\times k}$ randomly, then symmetric decorrelate
\For{$t=1$ to $\texttt{max\_iter}$}
  \State $\mathbf{Y}\gets \mathbf{W}\mathbf{Z}$
  \State $\mathbf{G}\gets \tanh(a\mathbf{Y})$, $\mathbf{G}'\gets a\,(1-\tanh^2(a\mathbf{Y}))$
  \State $\mathbf{D}\gets \operatorname{diag}(\mathrm{mean}(\mathbf{G}',\text{dims}=2))$
  \State $\mathbf{W}_{\text{new}}\gets \frac{1}{N}\mathbf{G}\mathbf{Z}^{\top}-\mathbf{D}\mathbf{W}$
  \State Symmetric decorrelation of $\mathbf{W}_{\text{new}}$
  \State $\Delta\gets \max_i\left|1-\left|\left(\mathbf{W}_{\text{new}}\mathbf{W}^{\top}\right)_{ii}\right|\right|$
  \If{$\Delta<\tau$}
    \State $\mathbf{W}\gets \mathbf{W}_{\text{new}}$; \textbf{break}
  \EndIf
  \State $\mathbf{W}\gets \mathbf{W}_{\text{new}}$
\EndFor
\State $\mathbf{S}\gets \mathbf{W}\mathbf{Z}$, $\mathbf{W}_{\text{total}}\gets \mathbf{W}\mathbf{V}_{\text{whit}}$, $\mathbf{A}\gets \mathbf{W}_{\text{total}}^{-1}$
\Ensure $\mathbf{S}$, $\mathbf{W}_{\text{total}}$, $\mathbf{A}$
\end{algorithmic}
\end{algorithm}

In Phase 2b, each IC is evaluated to distinguish neural from artifactual activity before reconstruction. The stage takes as input $(\mathbf{S},\mathbf{W}_{\text{total}},\mathbf{A})$, channel metadata, and optionally electrode coordinates for topographic inspection.

Spatial, spectral, and statistical features are standardized (z-score) and combined into four artifact-oriented scores (ocular, muscle, line, jump). The global decision index is defined as
\sidenote{%
\scriptsize
\textbf{ocular\_score}\\
Captura patrones de parpadeo/EOG:\\
predominio frontal, energia lenta\\
y, cuando existe, correlacion con EOG.
}
\sidenote{%
\scriptsize
\textbf{muscle\_score}\\
Resume evidencia de EMG:\\
potencia alta en bandas rapidas\\
y trazas de alta frecuencia.
}
\[
\begin{aligned}
\mathrm{artifact\_score}_i
&=
\max\!\big(
\mathrm{ocular\_score}_i,\,
\mathrm{muscle\_score}_i,\\
&\qquad
\mathrm{line\_score}_i,\,
\mathrm{jump\_score}_i
\big).
\end{aligned}
\]
\sidenote{%
\scriptsize
\textbf{line\_score}\\
Indica contaminacion de red\\
(50/60 Hz y armonicos) en el IC.
}
\sidenote{%
\scriptsize
\textbf{jump\_score}\\
Refleja transientes abruptos\\
(saltos, picos, clipping parcial),\\
normalmente no fisiologicos.
}
Thus, each component is judged by its most suspicious feature family. ICs above threshold are labelled as artifact and removed by zeroing the corresponding rows in $\mathbf{S}$.
\sidenote{%
\scriptsize
\textbf{Interpretacion del max}\\
Si un solo tipo de artefacto es alto,\\
el IC se marca como sospechoso\\
aunque los demas scores sean bajos.
}

Clean EEG is then reconstructed by
\[
\mathbf{X}_{\text{clean}}=\mathbf{A}\mathbf{S}_{\text{clean}}.
\]
This equation restores the signal in channel space using only retained components. Final outputs include the clean channel-indexed EEG and full audit metadata (bad IC list, feature table, labels, and optional comparison plots) to preserve reproducibility and traceability.
\sidenote{%
\scriptsize
\textbf{Umbral practico}\\
Regla tipica: \texttt{artifact\_score > 1.5}.\\
Mas bajo: mas sensible;\\
mas alto: mas conservador.
}


\subsection{Phases 3--6 (Segment processing)}

After ICA cleaning (Phase 2b), continuous EEG is segmented into fixed-length epochs to build a stable representation for spectral and connectivity analysis. The segmentation stage converts each channel time series into a 3D tensor
\[
\mathbf{E}\in\mathbb{R}^{C\times L\times K},
\]
where $C$ is the number of channels, $L$ is samples per epoch, and $K$ is the number of valid epochs.
\begin{figure}[H]
\centering
\resizebox{0.62\columnwidth}{!}{%
\begin{tikzpicture}[>=Latex]
  % Front face
  \draw[rounded corners=1.5pt, draw=black!75, fill=black!3] (0,0) rectangle (3.2,2.1);
  % Back face offset
  \draw[draw=black!55, fill=black!8] (0.9,0.7) rectangle (4.1,2.8);
  % Connect corners
  \draw[black!55] (0,0)--(0.9,0.7);
  \draw[black!55] (3.2,0)--(4.1,0.7);
  \draw[black!55] (0,2.1)--(0.9,2.8);
  \draw[black!55] (3.2,2.1)--(4.1,2.8);
  % Inner hint lines
  \draw[black!25] (1.6,0)--(1.6,2.1);
  \draw[black!25] (0,1.05)--(3.2,1.05);
  \node at (2.15,1.55) {$\mathbf{E}$};
  % Axis arrows
  \draw[-Latex,thick] (-0.55,0)--(-0.55,2.1);
  \node[anchor=east] at (-0.62,1.05) {$C$};
  \draw[-Latex,thick] (0,-0.45)--(3.2,-0.45);
  \node[below] at (1.6,-0.45) {$L$};
  \draw[-Latex,thick] (3.2,2.1)--(4.4,3.05);
  \node[anchor=west] at (4.45,3.05) {$K$};
\end{tikzpicture}
}
\caption{Epoch tensor representation ($C \times L \times K$).}
\end{figure}

Epoch length is defined from sampling rate and target duration:
\[
L=\left\lfloor F_s\cdot T_{\mathrm{seg}}\right\rfloor.
\]
For example, with $T_{\mathrm{seg}}=1\,\mathrm{s}$ and $F_s=500\,\mathrm{Hz}$, each epoch contains $L=500$ samples. The start-to-start distance between consecutive epochs is
\[
\mathrm{step}=L-\mathrm{overlap}.
\]
With zero overlap, $\mathrm{step}=L$, so epochs are consecutive and statistically independent.

Given a signal with $N_{\mathrm{total}}$ samples, the number of complete epochs is
\[
K=\left\lfloor\frac{N_{\mathrm{total}}-L}{\mathrm{step}}\right\rfloor+1.
\]
Samples at the end that cannot fill a complete epoch are discarded. This avoids variable-length windows and keeps the tensor shape consistent across channels and subjects.

The 1-second default is a practical compromise: it gives spectral resolution
\[
\Delta f\approx\frac{F_s}{L}\approx 1\,\mathrm{Hz},
\]
which is adequate for band-based connectivity, while still supporting the weak-stationarity assumption within each epoch.

Operationally, \texttt{src/Processing/segmentation.jl} loads cleaned EEG from \texttt{data/Processing/ICA/}, performs the split, computes per-epoch descriptive statistics (mean, standard deviation, minimum, maximum, RMS), produces example plots (overlay and stacked epochs), and saves both tensors and metadata under \texttt{data/Processing/segmentation/}. These outputs are then baseline-corrected and passed to amplitude-based artifact rejection, with optional second baseline correction after rejection.

\paragraph{Baseline correction (before artifacts).}
After segmentation, each epoch may still contain a DC offset or slow local drift that differs across epochs. The first baseline correction step removes this offset by subtracting, from each epoch, the mean value within a predefined initial window. If the baseline window in seconds is $[t_{\mathrm{start}},t_{\mathrm{end}}]$, its sample indices are
\[
n_{\mathrm{start}}=\lfloor t_{\mathrm{start}}F_s\rfloor+1,\qquad
n_{\mathrm{end}}=\lfloor t_{\mathrm{end}}F_s\rfloor.
\]
For channel $c$ and epoch $k$, with $N_{\mathrm{bl}}=n_{\mathrm{end}}-n_{\mathrm{start}}+1$, the baseline mean is
\[
\mu_{c,k}=\frac{1}{N_{\mathrm{bl}}}\sum_{n=n_{\mathrm{start}}}^{n_{\mathrm{end}}}E_{c,n,k},
\]
and the corrected epoch is
\[
E^{\mathrm{corr}}_{c,:,k}=E_{c,:,k}-\mu_{c,k}.
\]
This keeps within-epoch shape and variance unchanged while setting the baseline window around zero, reducing bias in amplitude and spectral metrics. In this project, the default window is the first 100\,ms ($t_{\mathrm{start}}=0$, $t_{\mathrm{end}}=0.1$\,s), which at $F_s=500$\,Hz corresponds to samples 1--50.
\begin{figure}[H]
\centering
\resizebox{0.86\columnwidth}{!}{%
\begin{tikzpicture}[x=1cm,y=1cm,>=Latex,font=\scriptsize]
  % Axes
  \draw[->, gray!70] (-2.8,-1.8) -- (4.5,-1.8)
    node[right] {time relative to event (ms)};
  \draw[->, gray!70] (-2.8,-1.6) -- (-2.8,1.4)
    node[above, rotate=90, anchor=south] {EEG amplitude ($\mu$V)};

  % Time ticks
  \draw[gray!65] (-2,-1.86)--(-2,-1.74);
  \draw[gray!65] (0,-1.86)--(0,-1.74);
  \node[below] at (-2,-1.8) {$-200$};
  \node[below] at (0,-1.8) {$0$};

  % Event marker
  \draw[dashed, gray] (0,-1.8)--(0,1.15);
  \node[anchor=west] at (0.08,1.08) {event / stimulus onset};

  % Baseline window region
  \fill[blue!8] (-2, -1.25) rectangle (0, 1.05);
  \draw[blue!35, dashed] (-2,-1.25) rectangle (0,1.05);
  \node[anchor=west] at (-1.9,0.95) {baseline window};

  % Shared waveform shape and corrected offset
  \draw[blue!75!black, line width=1.05pt, domain=-2.4:4.0, samples=120]
    plot (\x,{0.28 + 0.23*sin(95*\x) + 0.10*cos(185*\x)});
  \draw[teal!75!black, line width=1.05pt, domain=-2.4:4.0, samples=120]
    plot (\x,{-0.62 + 0.23*sin(95*\x) + 0.10*cos(185*\x)});

  % Baseline mean in pre-event window
  \draw[dashed, gray!70] (-2,0.28)--(0,0.28);
  \node[anchor=east] at (-2.05,0.28) {$\mu_{\mathrm{baseline}}$};
  \node[anchor=west, gray!70] at (-1.95,0.12) {mean over baseline window};

  % Gentle horizontal guides for levels
  \draw[gray!30] (-2.6,0.0)--(4.2,0.0);
  \draw[gray!30] (-2.6,-0.90)--(4.2,-0.90);

  % Correction arrow
  \draw[-Latex, very thick, black!70] (4.25,0.28)--(4.25,-0.62);
  \node[anchor=west] at (4.33,-0.17) {subtract $\mu_{\mathrm{baseline}}$};

  % Curve labels
  \node[anchor=west] at (2.35,0.78) {before correction};
  \node[anchor=west] at (2.35,-0.14) {after correction};

  % Formula in clear area
  \node[anchor=west] at (0.55,-1.35)
    {$x_{\mathrm{corrected}}(t)=x(t)-\mu_{\mathrm{baseline}}$};
\end{tikzpicture}
}
\caption{Conceptual effect of EEG/ERP baseline correction in a single epoch: a constant vertical shift using the pre-event baseline mean.}
\end{figure}

\paragraph{Why baseline both before and after artifact rejection?}
The first baseline correction (before artifact rejection) standardizes offsets so amplitude-threshold rules are comparable across epochs and channels. Without this step, epochs with different DC levels could be over- or under-rejected for purely offset-related reasons.

After artifact rejection, the retained epoch set is no longer identical to the pre-rejection set, and truncated/removed epochs can shift dataset-level central tendency. An optional second baseline correction re-centers the final cleaned tensor used for FFT/CSD/wPLI, improving consistency across subjects and reruns. In short: the first baseline supports fair artifact decisions; the second baseline harmonizes the final analysis input.

\paragraph{Artifact rejection (amplitude-based).}
Even after ICA and baseline correction, some epochs can still include residual bursts (e.g., EMG spikes, electrode pops, brief saturations). The rejection step in \texttt{src/Processing/artifact\_rejection.jl} applies a conservative amplitude rule and removes whole epochs that violate configured bounds.

Let $\mathbf{E}^{\mathrm{bl}}\in\mathbb{R}^{C\times L\times K}$ be baseline-corrected epochs and let
\[
\mathcal{C}_{\mathrm{used}}=\{1,\dots,\min(N_{\mathrm{used}},C)\}
\]
be the channels inspected for thresholding (default $N_{\mathrm{used}}=30$). For each epoch $k$, define the global minimum and maximum over inspected channels and all samples:
\[
m_k=\min_{c\in\mathcal{C}_{\mathrm{used}},\,1\le n\le L}E^{\mathrm{bl}}_{c,n,k},
\qquad
M_k=\max_{c\in\mathcal{C}_{\mathrm{used}},\,1\le n\le L}E^{\mathrm{bl}}_{c,n,k}.
\]
Epoch $k$ is marked as bad if
\[
m_k<A_{\min}\quad\text{or}\quad M_k>A_{\max},
\]
with default conservative limits $A_{\min}=-70\,\mu\mathrm{V}$ and $A_{\max}=+70\,\mu\mathrm{V}$.
\sidenote{%
\footnotesize
\textbf{Threshold logic (single epoch)}\\[1pt]
\centering
\begin{tikzpicture}[x=0.95cm,y=0.55cm,>=Latex]
  \draw[black!35] (0,0)--(4.8,0);
  \draw[red!70!black, dashed, line width=0.7pt] (0,0.85)--(4.8,0.85);
  \draw[red!70!black, dashed, line width=0.7pt] (0,-0.85)--(4.8,-0.85);
  \node[anchor=west, scale=0.68] at (4.85,0.85) {$A_{\max}$};
  \node[anchor=west, scale=0.68] at (4.85,-0.85) {$A_{\min}$};
  \draw[blue!75!black, line width=1.0pt]
    (0.1,0.10) .. controls (1.1,0.45) and (2.2,-0.20) .. (3.0,0.25)
              .. controls (3.5,0.55) and (3.8,1.10) .. (4.25,0.65);
  \node[scale=0.70] at (2.0,1.18) {violation};
\end{tikzpicture}\\
\textit{reject epoch if any sample crosses bounds}
}

The retained epochs form the filtered tensor $\mathbf{E}^{\mathrm{AR}}$ by concatenating only good indices. A useful quality metric is the rejection rate
\[
\mathrm{rate}=\frac{K-K_{\mathrm{AR}}}{K},
\]
where $K_{\mathrm{AR}}$ is the number of retained epochs. For diagnostics, the violating-channel set
\[
\mathcal{V}_k=
\left\{
c\in\mathcal{C}_{\mathrm{used}}:\,
\min_n E^{\mathrm{bl}}_{c,n,k}<A_{\min}
\;\text{or}\;
\max_n E^{\mathrm{bl}}_{c,n,k}>A_{\max}
\right\}
\]
is saved together with per-epoch min/max statistics.
\sidenote{%
\footnotesize
\textbf{Default processing values}\\
Segment length: \texttt{1 s}\\
Overlap: \texttt{0}\\
AR thresholds: \texttt{$\pm$70 $\mu$V}\\
Used channels for AR: \texttt{30}\\
Baseline window: \texttt{0--0.1 s}\\
Second baseline: optional
}

\subsection{Phase 7 (Connectivity)}

\subsubsection*{7.1 Current Source Density (CSD) transformation}
Volume conduction spreads currents through head tissues, so each scalp electrode records a weighted mixture of multiple sources. This inflates zero-lag coupling between nearby channels and reduces spatial specificity. To improve interpretability before connectivity estimation, the pipeline applies a Current Source Density (CSD) transform, i.e., a surface Laplacian approximation of the second spatial derivative of scalp potential \citep{Cohen2014}.

Operationally, \texttt{src/Connectivity/CSD.jl} loads second-baseline-corrected segments and electrode coordinates (BIDS TSV), checks channel-electrode consistency, and applies a Perrin-style spherical spline Laplacian (BVA-equivalent). Let
\[
\mathbf{E}\in\mathbb{R}^{C\times L\times K}
\]
be the segmented tensor and $\mathbf{L}\in\mathbb{R}^{C\times C}$ the CSD operator. For each segment $s$,
\[
\mathbf{E}_{\mathrm{CSD}}[:,:,s]=\mathbf{L}\,\mathbf{E}[:,:,s].
\]
This yields an effectively reference-free representation (constant component removed by construction), better suited for phase-based connectivity.

Electrode coordinates are first normalized to the unit sphere:
\[
r_c=\sqrt{x_c^2+y_c^2+z_c^2},\qquad
\mathbf{R}_c=\left(\frac{x_c}{r_c},\frac{y_c}{r_c},\frac{z_c}{r_c}\right).
\]
\begin{figure}[H]
\centering
\resizebox{0.78\columnwidth}{!}{%
\begin{tikzpicture}[font=\scriptsize]
  % Head outline
  \draw[thick] (0,0) circle (2.6);
  % Nose
  \draw[thick] (-0.35,2.58) -- (0,2.95) -- (0.35,2.58);
  % Ears
  \draw[thick] (-2.6,0.65) .. controls (-2.95,0.25) and (-2.95,-0.25) .. (-2.6,-0.65);
  \draw[thick] (2.6,0.65) .. controls (2.95,0.25) and (2.95,-0.25) .. (2.6,-0.65);
  % Axes
  \draw[-Latex,black!70] (-3.2,-3.0) -- (-1.8,-3.0) node[below] {$x$};
  \draw[-Latex,black!70] (-3.2,-3.0) -- (-3.2,-1.6) node[left] {$y$};

  % Electrode positions (2D approximation in normalized head coordinates)
  \foreach \name/\x/\y in {
    Fp1/-0.95/2.05, Fp2/0.95/2.05,
    F7/-2.00/1.20, F3/-0.95/1.25, Fz/0/1.35, F4/0.95/1.25, F8/2.00/1.20,
    T7/-2.25/0.15, C3/-1.05/0.20, Cz/0/0.25, C4/1.05/0.20, T8/2.25/0.15,
    P7/-1.90/-1.00, P3/-0.95/-0.85, Pz/0/-0.90, P4/0.95/-0.85, P8/1.90/-1.00,
    O1/-0.75/-2.00, Oz/0/-2.10, O2/0.75/-2.00,
    TP9/-2.25/-1.55, TP10/2.25/-1.55
  }{
    \filldraw[blue!65!black] (\x,\y) circle (1.2pt);
    \node[black] at (\x,\y+0.19) {\name};
  }
\end{tikzpicture}
}
\caption{Approximate 2D scalp layout (x-y view) with labeled electrodes used for spatial operators such as CSD.}
\end{figure}

Using dot products $\cos\gamma_{ij}=\mathbf{R}_i\cdot\mathbf{R}_j$, spline kernels are built through truncated Legendre expansions (up to \texttt{degree}):
\[
G_{ij}=\sum_{k=1}^{\texttt{degree}} c_g(k)\,P_k(\cos\gamma_{ij}),
\qquad
H_{ij}=\sum_{k=1}^{\texttt{degree}} c_h(k)\,P_k(\cos\gamma_{ij}).
\]
After regularization $G_\lambda=G+\lambda I$, the reference-free constraint is enforced with
\[
 C
 =
 G_\lambda^{-1}
 -
 \frac{(G_\lambda^{-1}\mathbf{1})(\mathbf{1}^{\top}G_\lambda^{-1})}
      {\mathbf{1}^{\top}G_\lambda^{-1}\mathbf{1}},
\]
and the final CSD operator is
\[
\mathbf{L}=H\,C.
\]
\sidenote{%
\scriptsize
\textbf{CSD defaults (BVA-aligned)}\\
$m$ (spline order): \texttt{4}\\
\texttt{degree}: \texttt{10}\\
$\lambda$: $10^{-5}$\\
Script: \texttt{src/Connectivity/CSD.jl}
}

\begin{algorithm}[H]
\caption{CSD via Perrin spherical spline Laplacian (as in \texttt{CSD.jl})}
\begin{algorithmic}
\Require $\mathbf{E}\in\mathbb{R}^{C\times L\times K}$, BIDS electrode TSV $(x_c,y_c,z_c)$, $m$, \texttt{degree}, $\lambda$
\State Load EEG-type electrode rows and match order to channel list
\If{any channel lacks coordinates}
  \State abort with consistency error
\EndIf
\For{$c=1$ to $C$}
  \State Normalize position to unit sphere: $\mathbf{R}_c\gets (x_c,y_c,z_c)/\sqrt{x_c^2+y_c^2+z_c^2}$
\EndFor
\State Precompute Perrin coefficients $c_g(k), c_h(k)$ for $k=1,\dots,\texttt{degree}$
\For{$i=1$ to $C$}
  \For{$j=i+1$ to $C$}
    \State $\cos\gamma \gets \mathrm{clamp}(\mathbf{R}_i\!\cdot\!\mathbf{R}_j,\,-1,1)$
    \State Build $G_{ij}$ and $H_{ij}$ via Legendre sums
  \EndFor
\EndFor
\State $G_\lambda\gets G+\lambda I$; enforce reference-free constraint to compute $C$
\State $\mathbf{L}\gets H\,C$
\For{$s=1$ to $K$}
  \State $\mathbf{E}_{\mathrm{CSD}}[:,:,s]\gets \mathbf{L}\,\mathbf{E}[:,:,s]$
\EndFor
\Ensure $\mathbf{E}_{\mathrm{CSD}}$ and exported CSD artifacts (matrices/figures/tables)
\end{algorithmic}
\end{algorithm}

\subsubsection*{7.2 Spectral decomposition and band-limited signals}
After optional second baseline correction (Phase 6), segmented EEG is used to estimate spectra in the frequency domain and derive band-wise power descriptors for downstream connectivity and quality analysis. The spectral routine (\texttt{src/Spectral/FFT.jl}) follows a BVA-compatible sequence on the segmented tensor:
\[
\mathbf{E}^{\mathrm{bl2}}\in\mathbb{R}^{C\times L\times K}.
\]

\begin{figure}[H]
\centering
\resizebox{0.98\columnwidth}{!}{%
\begin{tikzpicture}[
  node distance=4mm and 4mm,
  box/.style={draw=black!70, rounded corners=2pt, fill=black!3, align=center, inner sep=4pt, font=\scriptsize},
  arr/.style={-Latex, thick, black!75}
]
\node[box] (in) {$\mathbf{E}^{\mathrm{bl2}}$\\segment input};
\node[box, right=of in] (dc) {DC removal\\$E^{\mathrm{dc}}=E-\mu$};
\node[box, right=of dc] (win) {Windowing\\$E^{\mathrm{win}}=E^{\mathrm{dc}}\odot w$};
\node[box, right=of win] (pad) {Zero-padding\\$N_{\mathrm{fft}}=512$};
\node[box, right=of pad] (fft) {rFFT\\$X=\mathrm{rFFT}(E^{\mathrm{pad}})$};
\node[box, below=of fft] (pow) {Power spectrum\\$P_\ell=|X_\ell|^2$};
\node[box, left=of pow] (corr) {Corrections\\$\overline{w^2}$, fold, $N_{\mathrm{fft}}^2$};
\node[box, left=of corr] (bands) {Band averages\\$\delta,\theta,\alpha,\beta,\gamma$};

\draw[arr] (in) -- (dc);
\draw[arr] (dc) -- (win);
\draw[arr] (win) -- (pad);
\draw[arr] (pad) -- (fft);
\draw[arr] (fft) -- (pow);
\draw[arr] (pow) -- (corr);
\draw[arr] (corr) -- (bands);
\end{tikzpicture}
}
\caption{Conceptual flow used in Section 7.2 for FFT-based spectral decomposition and band-power estimation.}
\end{figure}

For each channel $c$ and segment $k$, a per-segment DC estimate is removed before tapering,
\[
\mu_{c,k}=\frac{1}{L}\sum_{n=1}^{L}E^{\mathrm{bl2}}_{c,n,k},
\qquad
E^{\mathrm{dc}}_{c,:,k}=E^{\mathrm{bl2}}_{c,:,k}-\mu_{c,k},
\]
then a Hamming/Tukey-style window $w[n]$ is applied:
\[
E^{\mathrm{win}}_{c,:,k}=E^{\mathrm{dc}}_{c,:,k}\odot w.
\]
Zero-padding to $N_{\mathrm{fft}}=512$ produces the FFT-ready tensor $\mathbf{E}^{\mathrm{pad}}$, giving frequency resolution
\[
\Delta f=\frac{F_s}{N_{\mathrm{fft}}}\approx 0.976\ \mathrm{Hz}\quad (F_s=500\ \mathrm{Hz}).
\]

The real FFT is computed per channel/segment:
\[
X_{c,:,k}=\mathrm{rFFT}\!\left(E^{\mathrm{pad}}_{c,:,k}\right),
\]
with $n_{\mathrm{bins}}=N_{\mathrm{fft}}/2+1$. DC is reintroduced in bin 0 (BVA convention), and power per bin is obtained from
\[
P_{c,\ell,k}=|X_{c,\ell,k}|^2,
\]
followed by window-variance correction (division by $\overline{w^2}$), full-spectrum folding (doubling interior bins), and normalization by $N_{\mathrm{fft}}^2$, so the final spectral power is expressed in $\mu V^2$.

Band-specific power is then computed by averaging bins whose center frequency lies inside each band interval $[f_{\min},f_{\max})$ (Delta, Theta, Alpha, Beta low/mid/high, Gamma), first per segment and then across segments.
\sidenote{%
\scriptsize
\textbf{FFT defaults}\\
$N_{\mathrm{fft}}=512$\\
$\Delta f\approx 0.976$ Hz\\
Window: Hamming/Tukey\\
Output unit: $\mu V^2$
}

\begin{algorithm}[H]
\caption{Spectral analysis via FFT (as in \texttt{FFT.jl})}
\begin{algorithmic}
\Require $\mathbf{E}^{\mathrm{bl2}}\in\mathbb{R}^{C\times L\times K}$, $F_s$, $N_{\mathrm{fft}}=512$, window parameters
\For{$c=1$ to $C$}
  \For{$k=1$ to $K$}
    \State $\mu_{c,k}\gets \mathrm{mean}(\mathbf{E}^{\mathrm{bl2}}[c,:,k])$
    \State $\mathbf{E}^{\mathrm{dc}}[c,:,k]\gets \mathbf{E}^{\mathrm{bl2}}[c,:,k]-\mu_{c,k}$
    \State Build taper $w[n]$ and apply $\mathbf{E}^{\mathrm{win}}[c,:,k]\gets \mathbf{E}^{\mathrm{dc}}[c,:,k]\odot w$
    \State Zero-pad to length $N_{\mathrm{fft}}$: $\mathbf{E}^{\mathrm{pad}}[c,1\!:\!L,k]\gets \mathbf{E}^{\mathrm{win}}[c,:,k]$
    \State $X_{c,:,k}\gets \mathrm{rFFT}(\mathbf{E}^{\mathrm{pad}}[c,:,k])$
    \State Reintroduce DC in bin 0 and compute power $P_{c,:,k}\gets |X_{c,:,k}|^2$
    \State Apply variance correction, spectrum folding, and normalization by $N_{\mathrm{fft}}^2$
  \EndFor
\EndFor
\State For each band $b$, average $P$ over bins with $f_\ell\in[f_{\min}^{(b)},f_{\max}^{(b)})$
\Ensure Power-per-bin, segment-averaged spectra, and band-wise power outputs
\end{algorithmic}
\end{algorithm}

\subsubsection*{7.3 wPLI connectivity metrics}
Phase-based functional connectivity is estimated with weighted Phase Lag Index (wPLI), which emphasizes consistent non-zero phase relations and downweights near-zero-lag contributions likely caused by volume conduction. For each frequency band, CSD signals are band-pass filtered and converted into analytic representations via Hilbert transform, then wPLI is aggregated across all samples and segments.

\begin{figure}[H]
\centering
\resizebox{0.98\columnwidth}{!}{%
\begin{tikzpicture}[
  node distance=4mm and 4mm,
  box/.style={draw=black!70, rounded corners=2pt, fill=black!3, align=center, inner sep=4pt, font=\scriptsize},
  arr/.style={-Latex, thick, black!75}
]
\node[box] (csd) {$\mathbf{E}_{\mathrm{CSD}}$\\band-specific input};
\node[box, right=of csd] (bp) {Band-pass\\$x_f=\mathrm{BP}(x,f_1,f_2)$};
\node[box, right=of bp] (hil) {Analytic signal\\$z=\mathrm{Hilbert}(x_f)$};
\node[box, right=of hil] (cross) {Cross term\\$S_{xy}=z_i\overline{z_j}$};
\node[box, below=of cross] (im) {Imaginary part\\$\mathrm{Im}_{ij}=\operatorname{Im}(S_{xy})$};
\node[box, left=of im] (agg) {Aggregation\\$\frac{|\sum \mathrm{Im}_{ij}|}{\sum |\mathrm{Im}_{ij}|+\epsilon}$};
\node[box, left=of agg] (wmat) {Output\\$\mathbf{W}\in[0,1]^{C\times C}$};

\draw[arr] (csd) -- (bp);
\draw[arr] (bp) -- (hil);
\draw[arr] (hil) -- (cross);
\draw[arr] (cross) -- (im);
\draw[arr] (im) -- (agg);
\draw[arr] (agg) -- (wmat);
\end{tikzpicture}
}
\caption{Conceptual flow used in Section 7.3 for band-limited wPLI connectivity estimation.}
\end{figure}

Let
\[
S_{xy}(t)=z_i(t)\,\overline{z_j(t)}
\]
be the cross-spectrum between analytic signals at channels $i$ and $j$, and define
\[
\mathrm{Im}_{ij}(t)=\operatorname{Im}\!\left(S_{xy}(t)\right)=\operatorname{Im}\!\left(z_i(t)\overline{z_j(t)}\right).
\]
Then wPLI is
\[
\mathrm{wPLI}_{ij}
=
\frac{\left|\mathbb{E}\left[\mathrm{Im}_{ij}(t)\right]\right|}
{\mathbb{E}\left[\left|\mathrm{Im}_{ij}(t)\right|\right]+\epsilon}.
\]
Only the imaginary component is used because zero-lag coupling tends to dominate the real part under volume conduction/common reference. Compared to plain PLI (which uses only signs), wPLI weights each sample by $|\mathrm{Im}_{ij}(t)|$, reducing sensitivity to weak noisy phase differences and improving robustness in finite samples \citep{Vinck2011WPLI}.

The connectivity matrix $\mathbf{W}\in\mathbb{R}^{C\times C}$ is symmetric, has diagonal zero, and values in $[0,1]$. Results are exported per band/condition as matrices, edge tables, heatmaps, and logs.
\sidenote{%
\scriptsize
\textbf{wPLI defaults}\\
Band-pass order: \texttt{8}\\
Filtering: \texttt{filtfilt} (zero-phase)\\
Aggregation: all $L\times K$ samples\\
$\epsilon$: \texttt{eps()}
}

\begin{algorithm}[H]
\caption{wPLI computation per band (as in \texttt{wPLI.jl})}
\begin{algorithmic}
\Require $\mathbf{E}_{\mathrm{CSD}}\in\mathbb{R}^{C\times L\times K}$, band $(f_1,f_2)$, $F_s$
\For{$c=1$ to $C$}
  \For{$s=1$ to $K$}
    \State $x\gets \mathbf{E}_{\mathrm{CSD}}[c,:,s]$
    \State $x_f\gets \mathrm{bandpass\_filtfilt}(x,F_s,f_1,f_2,\mathrm{order}=8)$
    \State $Z_c[:,s]\gets \mathrm{analytic\_signal}(x_f)$ \Comment{Hilbert via FFT}
  \EndFor
\EndFor
\For{$i=1$ to $C$}
  \For{$j=i+1$ to $C$}
    \State $\mathrm{Im}_{ij}\gets \operatorname{Im}(Z_i\odot \overline{Z_j})$
    \State $\mathrm{num}\gets |\sum \mathrm{Im}_{ij}|$;\; $\mathrm{den}\gets \sum |\mathrm{Im}_{ij}|+\epsilon$
    \State $W[i,j]\gets \mathrm{num}/\mathrm{den}$;\; $W[j,i]\gets W[i,j]$
  \EndFor
\EndFor
\Ensure $\mathbf{W}$ (symmetric, diagonal 0), plus matrix CSV, edges CSV, heatmap PNG
\end{algorithmic}
\end{algorithm}

\subsection{Phase 8 (Inference)}

After empirical connectivity matrices are obtained, the planned final stage builds surrogate-based null distributions and performs multiple-comparison-aware statistical inference. In the current repository state, this phase is fully documented in the Pluto export, but the executable source module \texttt{src/Surrogate/Surrogate.jl} is still a placeholder. The following specification therefore acts as the implementation target for the next code milestone.

\subsubsection{8.1 Surrogate data model}
To assess significance, surrogate EEG signals are generated independently per channel via phase randomization in the Fourier domain. For each channel, if $X(f)$ is the Fourier transform of the original signal, surrogate spectra are defined as
\[
X^{(s)}(f)=|X(f)|\,e^{\mathrm{i}\phi^{(s)}(f)},
\qquad
\phi^{(s)}(f)\sim\mathcal{U}(0,2\pi),
\]
with conjugate symmetry enforced so inverse FFT yields real-valued signals.

This construction preserves per-channel power spectral density,
\[
\left|X^{(s)}(f)\right|^2=\left|X(f)\right|^2,
\]
while destroying cross-channel phase relationships. Under this null, observed connectivity is attributable to chance phase alignment rather than genuine coupling.

Operationally, each surrogate realization repeats the same Phase 7 processing chain (band filtering, analytic signal extraction, and wPLI computation), producing surrogate connectivity matrices per band and condition.

\subsubsection{8.2 Statistical inference framework}
Let $\mathrm{wPLI}^{\mathrm{real}}_{ij}$ be the empirical edge value and $\mathrm{wPLI}^{(s)}_{ij}$ surrogate realizations for edge $(i,j)$. The empirical p-value is
\[
p_{ij}
=
\frac{1+\#\{\mathrm{wPLI}^{(s)}_{ij}\ge\mathrm{wPLI}^{\mathrm{real}}_{ij}\}}
{1+N_s},
\]
where $N_s$ is the number of surrogates. This plus-one form avoids zero p-values in finite surrogate samples and yields a discrete empirical null support.

P-values are then adjusted for multiplicity across edges using False Discovery Rate (FDR, Benjamini--Hochberg; default $\alpha=0.05$). The final interpretation is therefore based on edges that survive both empirical null comparison and multiplicity control.
\sidenote{%
\footnotesize
\textbf{Inference knobs}\\
\texttt{N\_s}: number of surrogates\\
Correction: FDR (default)\\
Trade-off: runtime vs p-value precision
}
